\addcontentsline{toc}{section}{СПИСОК ИСПОЛЬЗОВАННЫХ ИСТОЧНИКОВ}
% LTeX: enabled=false

\begin{thebibliography}{99}
	\bibitem{phillips} Phillips, B. Android Programming: The Big Nerd Ranch Guide / B. Phillips, C. Stewart, K. Marsicano. – Atlanta : Big Nerd Ranch Guides, 2022. – 700 с. – ISBN 978-0137645541. – Текст : непосредственный.
	\bibitem{deytel} Дейтел, П. Android для разработчиков / П. Дейтел, Х. Дейтел. – Санкт-Петербург : Питер, 2022. – 896 с. – ISBN 978-5-4461-1234-5. – Текст : непосредственный.
	\bibitem{zhukov} Жуков, А. Л. Разработка мобильных приложений на Kotlin / А. Л. Жуков. – Москва : ДМК Пресс, 2023. – 320 с. – ISBN 978-5-93700-123-4. – Текст : непосредственный.
	\bibitem{jemerov} Jemerov, D. Kotlin in Action / D. Jemerov, S. Isakova. – Shelter Island : Manning Publications, 2017. – 360 с. – ISBN 978-1617293290. – Текст : непосредственный.
	\bibitem{murphy} Murphy, M. L. The Busy Coder's Guide to Android Development / M. L. Murphy. – CommonsWare, 2021. – 2500 с. – ISBN 978-0-9816780-0-9. – Текст : непосредственный.
	\bibitem{vyazovik} Вязовик, Н. А. Программирование на Kotlin / Н. А. Вязовик. – Москва : ДМК Пресс, 2022. – 280 с. – ISBN 978-5-93700-156-2. – Текст : непосредственный.
	\bibitem{larman} Ларман, К. Применение UML 2.0 и шаблонов проектирования / К. Ларман ; пер. с англ. А. Ю. Шелестовой. – Москва : Вильямс, 2013. – 736 с. – ISBN 978-5-8459-1185-8. – Текст : непосредственный.
	\bibitem{buch} Буч, Г. Введение в UML от создателей языка / Г. Буч, Дж. Рамбо, И. Якобсон. – Москва : ДМК Пресс, 2015. – 496 с. – ISBN 978-5-94074-644-7. – Текст : непосредственный.
	\bibitem{fowler} Фаулер, М. UML. Основы / М. Фаулер. – Санкт-Петербург : Символ-Плюс, 2016. – 192 с. – ISBN 978-5-93286-056-4. – Текст : непосредственный.
	\bibitem{freeman} Фримен, Э. Паттерны проектирования / Э. Фримен, Э. Робсон, К. Сьерра. – Санкт-Петербург : Питер, 2018. – 656 с. – ISBN 978-5-496-03210-0. – Текст : непосредственный.
	\bibitem{martin} Мартин, Р. Чистая архитектура. Искусство разработки программного обеспечения / Р. Мартин ; пер. с англ. – Санкт-Петербург : Питер, 2018. – 351 с. – ISBN 978-5-4461-0772-8. – Текст : непосредственный.
	\bibitem{martinclean} Мартин, Р. Чистый код. Создание, анализ и рефакторинг / Р. Мартин. – Санкт-Петербург : Питер, 2019. – 464 с. – ISBN 978-5-4461-0960-9. – Текст : непосредственный.
	\bibitem{fowlerref} Фаулер, М. Рефакторинг. Улучшение существующего кода / М. Фаулер. – Санкт-Петербург : Питер, 2019. – 448 с. – ISBN 978-5-4461-1288-7. – Текст : непосредственный.
	\bibitem{somerville} Соммервилл, И. Инженерия программного обеспечения / И. Соммервилл. – Москва : Вильямс, 2002. – 624 с. – ISBN 5-8459-0330-0. – Текст : непосредственный.
	\bibitem{pressman} Прессман, Р. Разработка программного обеспечения: практическое руководство / Р. Прессман. – Москва : Вильямс, 2016. – 912 с. – ISBN 978-5-8459-2035-5. – Текст : непосредственный.
	\bibitem{vigers} Вигерс, К. Разработка требований к программному обеспечению / К. Вигерс, Д. Битти. – 3-е изд., доп. – Санкт-Петербург : BHV, 2020. – 736 с. – ISBN 978-5-9775-3348-5. – Текст : непосредственный.
	\bibitem{belik} Белик, А. Г. Проектирование и архитектура программных систем: учебное пособие / А. Г. Белик, В. Н. Цыганенко. – Омск : ОмГТУ, 2016. – 96 с. – ISBN 978-5-8149-2258-8. – Текст : непосредственный.
	\bibitem{gvozdeva} Гвоздева, Т. В. Проектирование информационных систем. Стандартизация: учебное пособие / Т. В. Гвоздева, Б. А. Баллод. – Санкт-Петербург : Лань, 2019. – 252 с. – ISBN 978-5-8114-3517-3. – Текст : непосредственный.
	\bibitem{gonzalez} Gonzalez, R. C. Digital Image Processing / R. C. Gonzalez, R. E. Woods. – 4th ed. – New York : Pearson, 2018. – 1024 с. – ISBN 978-0-13-335672-4. – Текст : непосредственный.
	\bibitem{sayood} Sayood, K. Introduction to Data Compression / K. Sayood. – 5th ed. – Cambridge : Morgan Kaufmann, 2017. – 790 с. – ISBN 978-0-12-809474-7. – Текст : непосредственный.
	\bibitem{salomon} Salomon, D. Data Compression: The Complete Reference / D. Salomon. – 4th ed. – London : Springer, 2007. – 1092 с. – ISBN 978-1-84628-602-5. – Текст : непосредственный.
	\bibitem{stevens} Стивенс, Р. UNIX: Профессиональное программирование / Р. Стивенс. – Санкт-Петербург : Питер, 2015. – 1024 с. – ISBN 978-5-496-01110-5. – Текст : непосредственный.
	\bibitem{tanenbaum} Таненбаум, Э. Современные операционные системы / Э. Таненбаум. – 4-е изд. – Санкт-Петербург : Питер, 2019. – 1120 с. – ISBN 978-5-496-01395-6. – Текст : непосредственный.
	\bibitem{kernighan} Керниган, Б. Язык программирования C / Б. Керниган, Д. Ритчи. – 2-е изд. – Санкт-Петербург : Невский Диалект, 2009. – 352 с. – ISBN 978-5-8459-0891-9. – Текст : непосредственный.
	\bibitem{gost-19-102} ГОСТ 19.102-77. Единая система программной документации. Стадии разработки. – Москва : Стандартинформ, 2010. – 12 с. – Текст : непосредственный.
	\bibitem{gost-34-601} ГОСТ 34.601-90. Информационная технология. Комплекс стандартов на автоматизированные системы. Стадии создания. – Москва : Стандартинформ, 2009. – 15 с. – Текст : непосредственный.
	\bibitem{gost-34-602} ГОСТ 34.602-89. Техническое задание на создание автоматизированной системы. – Москва : Стандартинформ, 2009. – 11 с. – Текст : непосредственный.
	\bibitem{gost-7-1} ГОСТ 7.1-2003. Библиографическая запись. Библиографическое описание. Общие требования и правила составления. – Москва : Стандартинформ, 2004. – 47 с. – Текст : непосредственный.
	\bibitem{zaytsev} Зайцев, А. А. Анализ мобильных приложений для фитнеса / А. А. Зайцев. – Текст : непосредственный // Молодой учёный. – 2025. – № 5. – С. 45–49.
	\bibitem{smirnova} Смирнова, Е. В. Использование SharedPreferences в Android / Е. В. Смирнова. – Текст : непосредственный // Информационные технологии. – 2024. – № 3. – С. 23–27.
	\bibitem{android-doc} Официальная документация Android Developers. – URL: https://developer.android.com/docs (дата обращения: 15.06.2026). – Текст : электронный.
	\bibitem{kotlin-doc} Официальная документация Kotlin. – URL: https://kotlinlang.org/docs/home.html (дата обращения: 15.06.2026). – Текст : электронный.
	\bibitem{sharedprefs} SharedPreferences. Документация Android Developers. – URL: https://developer.android.com/training/data-storage/shared-preferences (дата обращения: 15.06.2026). – Текст : электронный.
	\bibitem{mpandroidchart} PhilJay. MPAndroidChart. – URL: https://github.com/PhilJay/MPAndroidChart (дата обращения: 15.06.2026). – Текст : электронный.
	\bibitem{gramota} Правила русской орфографии и пунктуации. – URL: http://new.gramota.ru/spravka/rules (дата обращения: 15.06.2026). – Текст : электронный.
\end{thebibliography}